\documentclass[11pt,a4paper]{article}

\usepackage[utf8]{inputenc}
\usepackage[T1]{fontenc}
\usepackage[ngerman]{babel}
\usepackage[a4paper,left=22mm,right=22mm,top=18mm,bottom=20mm]{geometry}
\usepackage{xcolor}
\usepackage{hyperref}

% Siemens-inspirierte Farben und Schrift passend zum Lebenslauf.
\definecolor{CVPrimary}{HTML}{007A7A}
\definecolor{CVText}{HTML}{000028}
\definecolor{CVMuted}{HTML}{4B5B65}
\definecolor{CVRule}{HTML}{B8DCDC}
\hypersetup{colorlinks=true,urlcolor=CVPrimary}
\renewcommand{\familydefault}{\sfdefault}
\pagestyle{empty}
\setlength{\parindent}{0pt}
\setlength{\parskip}{0.65em}
\setlength{\emergencystretch}{2em}

% Für ein festes Versanddatum beispielsweise: 14. Oktober 2026
\newcommand{\Briefdatum}{\today}

\begin{document}
\color{CVText}

{\raggedleft \Briefdatum\par}
\vspace{5mm}
{\large\bfseries\textcolor{CVPrimary}{Bewerbung für das Siemens Graduate Program}}\par
{\small Digital Industries -- Automation, Tech Unit Motion | Job-ID 524054}\par
{\color{CVPrimary}\rule{\linewidth}{0.6pt}}\par
\vspace{1mm}

Sehr geehrte Damen und Herren,

Industrielle Automatisierung wird zunehmend von KI und Software geprägt. Genau an dieser Schnittstelle möchte ich im Siemens Graduate Program in der Tech Unit Motion arbeiten. Besonders reizt mich die Kombination aus industrieller Technologie, wechselnden Einsätzen und internationaler Zusammenarbeit, mit der Möglichkeit, verschiedene Bereiche kennenzulernen und eine Station im Ausland zu absolvieren.

Im dualen Bachelorstudium der Wirtschaftsinformatik bei Bosch durchlief ich verschiedene Abteilungen und absolvierte ein Praxissemester in Polen. Dabei lernte ich, mich zügig in neue Themen und Teams einzuarbeiten. Zudem präsentierte ich regelmäßig auch nicht-technische Themen anforderungsgerecht für fachfremde Zielgruppen. Die Arbeit in einem internationalen Umfeld hat mir besonders gefallen, und diese Erfahrung möchte ich im Graduate Program ausbauen.

Im Masterstudium Machine Learning \& Data Analytics an der Hochschule Aalen vertiefte ich meine Kenntnisse in Deep Learning, Natural Language Processing und Computer Vision. Parallel arbeitete ich als Werkstudentin bei ZEISS SMT an Optimierungsverfahren für das Job-Shop-Scheduling-Problem sowie an der Analyse und dem Monitoring der Maschinenverfügbarkeit. Dabei verband ich KI-Methoden mit konkreten industriellen Fragestellungen und stimmte mich eng mit den Fachabteilungen ab.

Bei Selfnet e.,V., einer studentischen Initiative zum Betrieb der Internetinfrastruktur in Studentenwohnheimen, wirke ich an der Entwicklung des Mitgliederverwaltungssystems mit und sammle praktische DevOps-Erfahrung. Als Mentor begleite ich neue Mitglieder, darunter viele internationale Studierende, beim Einstieg. Diese Teamarbeit und Eigeninitiative möchte ich in die Tech Unit Motion einbringen.

An der Tech Unit Motion reizt mich besonders, wie Industrial AI und Softwareentwicklung die Automatisierung weiterentwickeln. Ich möchte meine Erfahrung mit datengetriebenen Lösungen einbringen und lernen, wie aus technischen Ideen im Zusammenspiel von Entwicklung, Produktmanagement und weiteren Funktionen marktfähige Lösungen entstehen.

Ich freue mich darauf, Sie in einem persönlichen Gespräch von meiner Motivation zu überzeugen und mehr über die Stationen des Programms zu erfahren.

Mit freundlichen Grüßen

\vspace{6mm}
Natalie Böhm

\end{document}
